% ============================================================
% Capitolo 3 — Programmazione a modalità mutuamente esclusive
% (Programmazione a modalità mutuamente esclusive.pdf)
% ============================================================
\chapter{Programmazione a modalità mutuamente esclusive}
\label{cap:modalita-esclusive}

\begin{note}[Origine]
Appunti dell'esercitazione del 9 marzo 2026. Nota: il programma svolto in aula usa nomi differenti da quelli usati nelle slide.
\end{note}

\section{Idea generale}
Nella programmazione a \textbf{task} è utile individuare \textbf{porzioni di codice che vengono eseguite in maniera mutuamente esclusiva}. Si usano \textbf{variabili booleane} per creare:
\begin{itemize}
    \item porzioni mutuamente esclusive nel codice (\textbf{stati});
    \item segnali e condizioni che fanno saltare da una porzione all'altra (\textbf{eventi}).
\end{itemize}

Si usa il costrutto \texttt{IF THEN} associato a variabili booleane per creare tali porzioni. Una porzione contiene:
\begin{itemize}
    \item una o più \textbf{azioni};
    \item una o più \textbf{transizioni}.
\end{itemize}
Una porzione di codice modella una \textbf{modalità}, o meglio uno \textbf{stato}.

\section{Transizioni}
\begin{defn}[Transizione]
Una transizione:
\begin{itemize}
    \item è contenuta in una porzione di codice;
    \item porta in un'altra porzione di codice;
    \item è guidata da un \textbf{evento}, modellato attraverso un segnale booleano;
    \item è racchiusa in un \texttt{IF THEN};
    \item deve essere eseguita \textbf{una sola volta}: ciò si ottiene rendendo falso il segnale che modella l'evento alla conclusione dell'\texttt{IF} (\textbf{«consumare l'evento»}).
\end{itemize}
\end{defn}

Una transizione può avvenire anche in «reazione» a una \textbf{condizione} che diventa vera (senza eventi): ad esempio si salta alla modalità «indietro» se l'indice \texttt{i} ha raggiunto un certo limite.

\section{Esempio: due modalità con contatore}
Programma con due modalità in cui un contatore viene incrementato o decrementato. Si passa da una modalità all'altra se (in alternativa):
\begin{itemize}
    \item il contatore ha superato un limite negativo o positivo ($-1000$, $1000$);
    \item un evento \texttt{step} viene inviato attraverso un button (pulsante) nella parte grafica.
\end{itemize}

Dichiarazione delle variabili:
\begin{lstlisting}[basicstyle=\ttfamily\small]
PROGRAM PLC_PRG
VAR
    i : INT := 0;
    //--- Modalita' oppure STATI
    avanti : BOOL := TRUE;
    indietro : BOOL := FALSE;
    //--- Evento
    step : BOOL := FALSE;
END_VAR
\end{lstlisting}

Le variabili booleane \texttt{avanti} e \texttt{indietro} rappresentano gli \textbf{stati} (modalità) ed è garantito che siano sempre mutuamente esclusivi (uno \texttt{TRUE} e l'altro \texttt{FALSE}); \texttt{step} è l'\textbf{evento} esterno che forza il cambio di modalità indipendentemente dal valore del contatore.

\begin{intuition}[Perché questo pattern è importante]
Questo schema (stati = porzioni di codice mutuamente esclusive, transizioni guidate da eventi o condizioni) è la traduzione diretta in ST di una \textbf{macchina a stati}: è il ponte tra il formalismo a stati finiti (Capitoli~\ref{cap:formalismi-1} e~\ref{cap:formalismi-2}) e la programmazione concreta dei PLC, e sarà riutilizzato in tutti i casi di studio successivi (timer, stati temporizzati, funzioni di sicurezza).
\end{intuition}
